\section{IoT, Keamanan, dan Kebutuhan Kontrol Akses}

\textit{Internet of Things} (IoT) merupakan ekosistem yang menghubungkan berbagai perangkat fisik, sensor, aplikasi, dan layanan melalui jaringan komunikasi untuk mendukung pertukaran data serta otomatisasi layanan secara \textit{real-time}. Lingkungan IoT bersifat heterogen karena terdiri dari perangkat dengan kemampuan, protokol, dan sistem operasi yang berbeda-beda, serta terdistribusi karena perangkat tersebar pada berbagai lokasi dan saling berkomunikasi melalui jaringan yang luas. Selain itu, sebagian besar perangkat IoT memiliki keterbatasan sumber daya pada aspek komputasi, memori, energi, dan \textit{bandwidth}, sehingga mekanisme keamanan yang digunakan harus ringan dan efisien. Karakteristik IoT yang \textit{latency-sensitive} juga menuntut respons yang cepat karena banyak aplikasi IoT digunakan pada lingkungan kritis seperti kesehatan, industri, dan \textit{smart transportation} \parencite{kokila_authentication_2025}.

Dari sisi keamanan, lingkungan IoT memiliki \textit{attack surface} yang sangat luas akibat jumlah perangkat yang besar dan konektivitas yang terus meningkat. \textcite{williams_survey_2022} menjelaskan bahwa ancaman keamanan IoT dapat muncul pada berbagai lapisan, mulai dari perangkat, jaringan, hingga aplikasi. Kerentanan tersebut semakin diperparah oleh penggunaan \textit{firmware} yang tidak aman, mekanisme autentikasi yang lemah, serta pengelolaan \textit{trust} yang tidak memadai, sehiingga meningkatkan risiko akses tidak sah terhadap perangkat dan data IoT \parencite{altaibek_survey_2025}. Selain itu, konektivitas jaringan IoT yang sering berubah dan tidak stabil membuat proses pengamanan akses menjadi lebih kompleks karena kondisi perangkat dan konteks lingkungan dapat berubah secara dinamis \parencite{krishnasrija_dynamic_2026}.

Dalam konteks keeamanan IoT, autentikasi dan otorisasi memiliki fungsi yang berbeda tetapi saling berkaitan. Autentikasi bertujuan memverifikasi identitas pengguna atau perangkat, sedangkan otorisasi menentukan tindakan apa saja yang diperbolehkan setelah identitas tersebut berhasil diverifikasi \parencite{diaz_authorization_2025}. Pada lingkungan IoT yang dinamis, keputusan otorisasi tidak cukup hanya didasarkan pada identitas atau peran statis, tetapi juga perlu mempertimbangkan atribut kontekstual seperti lokasi, waktu, jenis perangkat, status keamanan perangkat, dan jenis aksi yang diminta. Oleh karena itu, IoT membutuhkan mekanisme kontrol akses yang \textit{fine-grained, context-aware, scalable, low-latency,} dan mampu beradaptasi terhadap perubahan konteks secara \textit{real-time}. \textcite{kalaria_adaptive_2024} menunjukkan bahwa \textit{adaptive context-aware access control} diperlukan agar kebijakan akses dapat menyesuaikan kondisi lingkungan secara dinamis sekaligus mengurangi ketergantungan pada \textit{cloud} melalui pemrosesan di \textit{fog} atau \textit{edge}. Sementara itu, \textcite{wang_edge-enabled_2025} menegaskan bahwa layananan IAM pada IoT masa depan harus mampu menyediakan \textit{access management} yang aman, skalabel, dan berlatensi rendah agar sesuai dengan kebutuhan layanan IoT.